
\documentclass{article}
\usepackage{subfig}
\usepackage{a4wide}
\usepackage[utf8]{inputenc}
\usepackage[russian]{babel}
\usepackage{graphicx}
\usepackage{amsmath}
\usepackage{float}
\usepackage{hyperref}
\usepackage{amssymb}
\usepackage{amsthm}
\usepackage{amsmath}
\usepackage{mathtools}
\usepackage{indentfirst}
\usepackage{hyperref}

\usepackage{setspace}
\usepackage{titlesec}
\setlength\parident{1.5em}
\begin{document}
\thispagestyle{empty}

\begin{center}
\ \vspace{-3cm}

\includegraphics[width=0.5\textwidth]{msu.eps}\\
{\scshape Московский государственный университет имени М.~В.~Ломоносова}\\
Факультет вычислительной математики и кибернетики\\
Кафедра системного анализа

\vfill

{\LARGE Отчет по практикуму}

\vspace{1cm}

{\Huge\bfseries <<Построение множества достижимости>>}
\end{center}

\vspace{1cm}

\begin{flushright}
  \large
  \textit{Студент 315 группы}\\
  Ю.\,М.~Алексеева

  \vspace{5mm}

  \textit{Руководитель практикума}\\
  м.н.с. И.\,А.~Чистяков
\end{flushright}

\vfill

\begin{center}
Москва, 2025
\end{center}

\newpage
\tableofcontents
\newpage
\section{Постановка задачи}
Задано обыкновенное дифференциальное уравнение:
$$\ddot{x} + \beta x -2 \sin(3x^3)+x\dot{x} = u,$$
где $x\in\mathbb{R}, ~ u\in\mathbb{R}.$ На возможные значения управляющего параметра $u$ наложено ограничение: $u \in [-\alpha, \alpha], \alpha > 0.$ Задан начальный момент времени $t_0 = 0$ и начальная позиция $x(t_0) = 0, ~\dot{x}(t_0) = 0.$ Необходимо построить множество достижимости $X(t, t_0, x(t_0), \dot{x}(t_0))$ (множество пар $(x(t), \dot{x}(t))$) в классе программных управлений в заданный момент времени $t\geqslant t_0$, а также исследовать его свойства.

\newpage
\section{Теоретические выкладки}
Введем замену $x_1 = x, x_2 = \dot{x}$, тогда уравнение сводится к системе вида:
$$
\begin{cases}
    \dot{x}_1 = x_2,\\
    \dot{x}_2 = -\beta x_1 + 2\sin(3x_1^3) - x_1x_2 + u.
\end{cases}
$$
\newtheorem{определение}{Определение}
\begin{определение}
Множество достижимости $X[t]$:
$$X[t] = X(t, t_0, x^0) =  \left\{x = x(t, u(\cdot)) \in \mathbb{R}^n :~\exists~  u(\cdot) , \text{ т.ч. }  u(\tau)\in \mathcal{P},~ \forall \tau\in[t_0, t],  ~x = x(t, t_0, x^0)\right\},$$
где $\mathcal{P}~-~$класс допустимых управлений (удовлетворяющих всем ограничениям).
\end{определение}
Введем функцию Гамильтона-Понтрягина для нашей задачи:
$$H(x, \psi, u) = \psi_1x_2 + \psi_2(-\beta x_1 + 2\sin(3x_1^3) - x_1x_2 + u).$$
\newtheorem{теорема}{Теорема}
\begin{теорема}[ПМП]\\
Пусть $(x^*(\cdot), ~u^*(\cdot))~-~$оптимальная пара на $[0, T],$ тогда $\exists~ \psi^*(t) \neq 0$ такая, что выполнено:
\begin{itemize}
    \item [1)] Сопряженная система (СС):
    $$
    \left.\dot{\psi}^*(t) = -\dfrac{\partial H}{\partial x}\right\vert_{\substack{\psi = \psi^*(t)\\x=x^*(t) \\ u=u^*(t)}};
    $$
    \item[2)] Условие максимума (УМ):
    $$
    H(x^*(t), \psi^*(t), u^*(t)) \overset{\text{п.в. t}}{=} \sup_{u\in\mathcal{P}}H(x^*(t), \psi^*(t), u);
    $$
    \item[3)] $\mathcal{M} (x^*(t), \psi^*(t)) = \text{const} \geqslant 0.$
\end{itemize}
\end{теорема}
Для максимизации функции Гамильтона-Понтрягина выбираем следующее управление:
$$
u^* = 
\begin{cases}
    \alpha, & \psi_2 > 0,\\
    [-\alpha, \alpha], & \psi_2 = 0,\\
    -\alpha, & \psi_2 < 0.
\end{cases}
$$
Возможно ли возникновение особого режима ($\psi_2(\tau) = 0$)?\\
Пусть $\exists ~\tau \in [0, T]: \psi_2(\tau) = 0$, тогда $\exists ~\delta>0: \forall t\in[\tau, \tau+\delta] ~ \psi_2(t) = 0, \dot{\psi}_2(t) = 0.$
Запишем сопряженную систему:
$$
\begin{cases}
    \dot{\psi}_1 = \beta \psi_2 - 18x_1^2\psi_2\cos(3x_1^3) + \psi_2 x_2,\\
    \dot{\psi}_2 = -\psi_1 + x_1\psi_2.
\end{cases}
$$
Отсюда видим, что в силу 2 уравнения сопряженной системы $\psi_1(t)= 0, ~\forall t\in[\tau, \tau+\delta].$ Таким образом, $\psi(t) = 0, \forall t\in[\tau, \tau+\delta],$ что противоречит условию ПМП о невырожденности $\psi(t).$ Значит, $u^* \overset{\text{п.в.}}{=} \alpha~\text{sign}(\psi_2).$\\
\begin{теорема}
Пусть $(x^*(\cdot), u^*(\cdot))~-~$оптимальная пара для времени быстродействия $T,$ удовлетворяющая ПМП, $\psi(t) = (\psi_1(t), \psi_2(t))~-~$соответствующая сопряженная функция. Тогда для любых $\tau_1, \tau_2$ таких, что $0<\tau_1<\tau_2<T$ выполнено:
\begin{itemize}
    \item [1.] если $\psi_2(\tau_1) = \psi_2(\tau_2) = 0$ и $x_2(\tau_1) = 0,$ то $x_2(\tau_2)= 0;$
    \item [2.] если $x_2(\tau_1) = x_2(\tau_2) = 0, x_2(t) \neq 0$ при $t\in (\tau_1, \tau_2)$ и $\psi_2(\tau_1) = 0,$ то $\psi_2(\tau_2) = 0;$
    \item [3.] если $\psi_2(\tau_1)=\psi_2(\tau_2)=0$ и $x_2(\tau_1)\neq 0,$ то $x_2(\tau_2)\neq 0$ и $\exists ~\tau'\in(\tau_1, \tau_2): ~x_2(\tau')= 0;$
    \item [4.] если $x_2(\tau_1)=x_2(\tau_2)=0, ~x_2(t)\neq 0$ при $t\in (\tau_1, \tau_2)$ и $\psi_2(\tau_1)\neq 0,$ то $\psi_2(\tau_2)\neq 0$ и $\exists ~\tau'\in(\tau_1, \tau_2): ~\psi_2(\tau')= 0.$
\end{itemize}
\end{теорема}
\renewcommand{\proofname}{Доказательство}
\begin{proof}
Напомним, что 
$$\mathcal{M} = \sup_{u\in[-\alpha, \alpha]}H(x, \psi, u) = \psi_1x_2+\psi_2(-\beta x_1 + 2\sin(3x_1^3) - x_1x_2) +\alpha\psi_2\text{sign}(\psi_2).$$
\begin{itemize}
    \item [1.] Т.к. $\mathcal{M} = \text{const},$ то $\mathcal{M}[t=\tau_1] = 0 = \psi_1(\tau_2)x_2(\tau_2).$ Т.к. $\psi_2(\tau_2) = 0,$ то в силу невырожденности вектора $\psi(t),$ получим, что $\psi_1(\tau_2) \neq 0 \implies x_2(\tau_2)=0.$
    \item [2.] Рассмотрим кусочно-непрерывную функцию $K(t) = \psi_1(t)x_2(t) + \psi_2(t)\dot{x}_2(t),$ определенная всюду, кроме моментов переключения. Пусть $\tau$ точка непрерывности $K,$ а $f = -\beta x_1+2\sin(3x_1^3)-x_1x_2,$ тогда при $t\in[\tau-\delta, \tau+\delta]:$
    $$\dfrac{dK}{dt} = \psi_2x_2\dfrac{\partial f}{\partial x_1}+\psi_1(-f + u) +(-\psi_1+x_1\psi_2)(-f+u)+\psi_2\dfrac{\partial^2 x_2}{\partial t^2}.$$
    При этом:
    $$\psi_2\dfrac{\partial^2 x_2}{\partial t^2} = -\psi_2\left(\dfrac{\partial f}{\partial x_1} x_2 + \dfrac{\partial f}{\partial x_2}(-f+u)\right).$$
    Объединяя полученные равенства, имеем:
    $$\dfrac{dK}{dt} = -\psi_2x_2\dfrac{\partial f}{\partial x_1}+\psi_1(-f + u) +(-\psi_1+x_1\psi_2)(-f+u)+ \psi_2\left(\dfrac{\partial f}{\partial x_1} x_2 - x_1(-f+u)\right)= 0.$$
    Пусть $t'~-~$момент переключения, тогда из ПМП $\psi_2(t')=0$,
    $$K(t'\pm 0) = \psi_1(t')x_2(t') \implies K(t'+0)=K(t'-0),$$
    т.е. $K(\cdot)$ непрерывна в точке $t'.$
    Доопределим $K$ во всех точках переключений и имеем, что $K(t) = \text{const}.$ Тогда ($\psi_2(\tau_1) = 0$):
    $$0 = \psi_2(\tau_1)\dfrac{dx_2(\tau_1\pm0)}{dt} = K(\tau_1) = K(\tau_2) = \psi_2(\tau_2)\dfrac{dx_2(\tau_2\pm0)}{dt}.$$
    От противного: \\
    Пусть $\psi_2(\tau_2) \neq 0,$ тогда $\dfrac{dx_2(\tau_2)}{dt} = 0,$ но по условию, $x_2(\tau_2) = \text{const}, u(t) = \text{const}$ при $t\in [\tau_2-\delta, \tau_2+\delta].$ Таким образом, $(x_1(\tau_2), x_2(\tau_2))~-~$стационарная точка системы
    $$
    \begin{cases}
        \dot{x}_1 = x_2,\\
        \dot{x}_2 = -f(x_1, x_2)+u.
    \end{cases}
    $$
    В силу единственности решения, имеем, что $x_1(t)= \text{const}$ при $t\in [\tau_2-\delta, \tau_2+\delta],$ что противоречит $x_2(t) \neq 0$ при $t\in(\tau_1, \tau_2).$ Значит, изначальное наше предположение неверно, значит, $\psi_2(\tau_2) = 0.$
    \item[3.] $\mathcal{M} = \text{const}.$ Без ограничения общности, считаем, что $\tau_1$ и $\tau_2~-~$два последовательных нуля $\psi_2(t)$, $\psi_2(t) \neq 0, ~\forall t\in(\tau_1, \tau_2).$ Так как $\dot{\psi}_2(\tau_i)=-\psi_1(\tau_i) \neq 0$ (из невырожденности вектора $\psi(t)$ и в силу того, что $\psi_2(\tau_i) = 0,$ следует, что $\psi_1(\tau_i)\neq0).$ Отсюда $\dot{\psi}_2(\tau_1)\dot{\psi}_2(\tau_2) < 0 \implies \psi_1(\tau_1)\psi_1(\tau_2) < 0.$ Тогда $\mathcal{M}[t = \tau_1] = \psi_1(\tau_1)x_2(\tau_1) = \psi_1(\tau_2)x_2(\tau_2) = \mathcal{M}[t = \tau_2] \implies x_2(\tau_1)x_2(\tau_2) < 0. $ Функция $x_2(t)$ непрерывная, $x_2(\tau_1)x_2(\tau_2) < 0 \implies \exists\tau': x_2(\tau') = 0.$
    \item[4.] По аналогии с пунктом 2, через противоречие со стационарной точкой получаем, что при $t\in(\tau_1, \tau_2):$
    $$(x_2(t))^2 + \left(\dfrac{dx_2(t\pm0)}{dt}\right)^2 \neq 0.$$
    Тогда
    $$0 \neq \psi_2(\tau_1) \dfrac{dx_2(\tau_1\pm0)}{dt} = K(\tau_1)=K(\tau_2) = \psi_2(\tau_2)\dfrac{dx_2(\tau_2\pm0)}{dt},$$
    то есть $\psi_2(\tau_2)\neq 0.$ Покажем, что $\exists t': \psi_2(t') = 0.$ 
    От противного:\\
    Пусть $\text{sign}(\psi_2(t))=\text{const}, \forall t \in(\tau_1,\tau_2).$ Тогда
    $$\dfrac{dx_2(\tau_1)}{dt}\dfrac{dx_2(\tau_2)}{dt}>0,$$
    это противоречит (по аналогии с 3 пунктом), что $\tau_1$ и $\tau_2~-~$последовательные нули $x_2.$
 \end{itemize}
 \end{proof}
Теорема доказана. Благодаря ей, делаем вывод, что нули $x_2(t), \psi_2(t)$ либо совпадают, либо чередуются.
\newpage
\section{Алгоритм решения}
Мы показали, что $u^* = \alpha\text{sign}(\psi_2).$ Тогда система приводится к виду:
$$
\begin{cases}
    \dot{x}_1 = x_2,\\
    \dot{x}_2 = -\beta x_1 + 2\sin(3x_1^3)-x_1x_2+\alpha\text{sign}(\psi_2).
\end{cases}
$$
Эту систему можно разбить на две системы $S^+, S^-$, которые отвечают значениям управления $u = \alpha, u = -\alpha$, соответственно:
$$
S^+ = 
\begin{cases}
    \dot{x}_1 = x_2,\\
    \dot{x}_2 = -\beta x_1 + 2\sin(3x_1^3)-x_1x_2+\alpha,
\end{cases}
S^- = 
\begin{cases}
    \dot{x}_1 = x_2,\\
    \dot{x}_2 = -\beta x_1 + 2\sin(3x_1^3)-x_1x_2-\alpha.
\end{cases}
$$
Начальное значение равно либо $u^* = \alpha, $ либ $u^* = -\alpha,$ рассмотрим два случая отдельно. Переключения управления могут происходить в моменты $\psi_2(\tau) = 0.$\\
Пусть управление $u^* = \alpha,$ Тогда решаем систему $S^+$ и находим момент времени $t': x(t')= 0.$ Строим полученную траекторию, соответствующую системе $S^+$ до момента времени $t'.$ Из теоремы о нулях мы знаем, что нули $\psi_2$ чередуются или совпадают с нулями $x_2.$ Тогда время переключения $\tau$ находится на отрезке $[0, t'].$ Сделаем перебор $\tau$ по этому отрезку. В момент переключения управление изменится на $u^* = -\alpha$ и начнет действовать система $S^-$. Тогда теперь из точки $(x_1(\tau), x_2(\tau))$ строить траекторию, соответствующую системе $S^-$ до нового момента переключения, который находится из решения сопряженной системы:
$$
\begin{cases}
    \dot{\psi}_1 = \beta \psi_2 - 18x_1^2\psi_2\cos(3x_1^3) + \psi_2 x_2,\\
    \dot{\psi}_2 = -\psi_1 + x_1\psi_2.
\end{cases}
$$
Найдем начальное условие для сопряженной системы. Возьмем момент переключения $\tau,$ в нем $\psi_2(\tau) = 0,$ а следовательно $\dot{\psi}_1 =0.$ Т.к. в начале $\psi_2>0,$ то $\dot{\psi}_2(\tau)<0,$ значит, $\psi_1(\tau) >0.$ Из положительной однородности функции Гамильтона-Понтрягина примем $\psi_1(\tau)=1.$ Тогда начальные условия для сопряженной системы:
$$
\begin{cases}
    \psi_1(\tau) = 1,\\
    \psi2(\tau)=0.
\end{cases}
$$
От момента переключения $\tau$ до следующего, которое мы находим из решения сопряженной системы, действует система $S^-,$ после $S^+$, и так до момента $t=T,$ момент останова.\\
Аналогичная процедура проводится для $u^* = -\alpha, \psi_2 < 0, \psi(\tau) = (-1, 0).$\\
\textbf{Алгоритм}:\\
Для построения множества достижимости в конкретный момент времени $T$ выразим концы каждой траектории, удовлетворяющей ПМП, а потом выберем из них только те точки, которые принадлежать границе множества достижимости.
\begin{itemize}
    \item [1.] Начинаем со случая $u^*= \alpha.$ Решаем систему $S^+,$ находим момент времени $x_2(t') = 0.$
    \item [2.] Запускаем перебор времени переключения $\tau$ по отрезку $[0, t'].$
    \item [3.] Решаем сопряженную систему с соответствующими начальными условиями.
    \item [4.] Делаем переключение и решаем систему $S^-$ с краевыми условиями $(x_1(\tau), x_2(\tau)),$ которые мы получили из прошлой системы.
    \item[5.] Решаем сопряженную систему с соответствующими для нее условиями (-1, 0).
    \item [6.] Делаем переключение и решаем систему $S^+.$
    \item [7.] Повторяем, пока $t < T.$
    \item [8.] Аналогично считаем случай, когда в начальный момент времени $u^*= -\alpha.$
    \item [9.] Объединяем полученную кривую в общее целое, которое является границей множества достижимости.
    \item [10.] При построении каждой траектории координаты точек переключения собираются в специальный массив, а затем объединяются в кривую переключений. 
\end{itemize}

\newpage
\section{Примеры работы программы}
На всех примерах розовым обозначена граница множества достижимости, синим$~-~$кривая переключений.
\subsection{Пример 1}
\begin{figure}[H]
    \centering
    \subfloat[$\alpha = 0.3, \beta = 1$]{
        \includegraphics[width=0.45\textwidth]{ex1.eps}
        \label{fig6}
    }
    \hfill
    \subfloat[$\alpha = 1, \beta = 0.3$]{
        \includegraphics[width=0.45\textwidth]{ex3.eps}
        \label{fig8}
    }
    \hfill
    \subfloat[$\alpha = 0.003, \beta = 0.001$]{
        \includegraphics[width=0.45\textwidth]{ex2.eps}
        \label{fig7}
    }
    \hfill
    \subfloat[$\alpha = 1, \beta = 3$]{
        \includegraphics[width=0.45\textwidth]{ex4.eps}
        \label{fig9}
    }
    \caption{Множество достижимости при $T = 1$ для различных $\alpha, \beta$}
     \label{fig69}
\end{figure}
\subsection{Пример 2}
\begin{figure}[H]
    \centering
    \subfloat[$T = 1$]{
        \includegraphics[width=0.45\textwidth]{ex5.eps}
        \label{fig6}
    }
    \hfill
    \subfloat[$T = 4$]{
        \includegraphics[width=0.45\textwidth]{ex6.eps}
        \label{fig8}
    }
    \hfill
    \subfloat[$T = 7$]{
        \includegraphics[width=0.45\textwidth]{ex7.eps}
        \label{fig7}
    }
    \hfill
    \subfloat[$T = 9.5$]{
        \includegraphics[width=0.45\textwidth]{ex8.eps}
        \label{fig9}
    }
    \caption{Множество достижимости при $\alpha = 0.01, \beta = 1$ для различных $T$}
     \label{fig69}
\end{figure}
\subsection{Пример 3}
\begin{figure}[H]
    \centering
    \subfloat[$\beta = 0.0001$]{
        \includegraphics[width=0.45\textwidth]{ex9.eps}
        \label{fig6}
    }
    \hfill
    \subfloat[$\beta = 0.01$]{
        \includegraphics[width=0.45\textwidth]{ex10.eps}
        \label{fig8}
    }
    \hfill
    \subfloat[$\beta = 1$]{
        \includegraphics[width=0.45\textwidth]{ex12.eps}
        \label{fig7}
    }
    \hfill
    \subfloat[$\beta = 10$]{
        \includegraphics[width=0.45\textwidth]{ex13.eps}
        \label{fig9}
    }
    \caption{Множество достижимости при $\alpha = 0.01, T = 0.9$ для различных $\beta$}
     \label{fig69}
\end{figure}


\newpage
\begin{thebibliography}{99}
\bibitem{1} Чистяков И.А., Паршиков М.В. Лекции по Оптимальному Управлению
\bibitem{2} Арутюнов А.В. Лекции по Многозначному Анализу.
\end{thebibliography}
\end{document}